\section{Resultados de verificación}
\label{sec:verificacion}

Cada resultado indica qué plataforma lo produjo. Ninguna medición eléctrica o de radio de esta sección
se hizo en la PCB, que no está fabricada.

\subsection{Sobre la PCB: verificación por diseño}

\begin{table}[H]
\centering
\footnotesize
\begin{tabularx}{\textwidth}{|l X X|}
\hline
\tablehead{\textbf{Verificación} & \textbf{Resultado} & \textbf{Evidencia}} \\ \hline
DRC de Altium & \bDrcRules{} reglas, \bDrcViolations{} violaciones (2026-11-15) & \ruta{evidence/drc_summary.txt} \\ \hline
Cálculos de potencia y carga & L, divisor MPPT, $R_{CS}$, TP4056, LDO y consumo (sección~3) & \ruta{docs/ARCHITECTURE.md} \\ \hline
Presupuesto de consumo en sueño & $\approx$ \cIsleep~\textmu A frente a \cIlimit~\textmu A, con pendientes & \ruta{hardware/POWER_BUDGET.md} \\ \hline
Salidas de fabricación & Gerbers, taladros, BOM y planos generados el 2026-11-18 & \ruta{hardware/pcb/} \\ \hline
Reporte ERC de Altium & 0 errores (informe del autor, 2026-11-26) & Informe de Altium, no incluido en el repositorio \\ \hline
\end{tabularx}
\caption{Verificaciones hechas sobre el diseño de la PCB.}
\end{table}

\subsection{Referencia indicativa: medidas en la \pBoard{}}

Estos números vienen de la prueba de concepto y valen \emph{solo} para esa plataforma y esas condiciones. No
describen la PCB, que usa otro procesador, otro receptor GNSS (con más constelaciones) y otra antena.

\begin{table}[H]
\centering
\footnotesize
\begin{tabularx}{\textwidth}{|X X|}
\hline
\tablehead{\textbf{Medida en la \pBoard{} (\pSoc)} & \textbf{Valor}} \\ \hline
\ruta{geofence_check()} con 4 vértices & \pBenchUsFour~\textmu s \\ \hline
\ruta{geofence_check()} con 1024 vértices & \pBenchMsThousand~ms; llega a 50~ms cerca de los \pBenchLimitVertices{} vértices \\ \hline
Dispersión horizontal, CEP50 (cielo parcial, menos de 15 min) & \pCepLow{} a \pCepHigh~m, frente a \pReqCep~m de RF-01 \\ \hline
Exactitud estimada por el receptor al aire libre, caminando & \pWalkAccMin{} a \pWalkAccMax~m (estimación del receptor, no un CEP medido) \\ \hline
Tiempo hasta el primer fix en frío & \pTtffA, \pTtffB{} y \pTtffC~s (cronómetro manual) \\ \hline
Memoria del firmware & \pFlashBytes~B de flash y \pRamBytes~B de RAM \\ \hline
\end{tabularx}
\caption{Resultados de la prueba de concepto, a título indicativo.}
\end{table}
